\chapter{Cholinesterase}

\section*{What Is This Test?}
Cholinesterase testing measures enzymes that break down acetylcholine. Plasma or serum pseudocholinesterase (butyrylcholinesterase) is used mainly in suspected organophosphate or carbamate poisoning and in selected assessments of prolonged paralysis after the anesthetic succinylcholine.

\section*{Alternative Names}
Pseudocholinesterase; Butyrylcholinesterase; Plasma Cholinesterase; Serum Cholinesterase.

\section*{Why This Test Is Ordered}
Testing can support diagnosis and monitoring of pesticide or nerve-agent exposure and evaluation of suspected inherited pseudocholinesterase deficiency after unexpectedly prolonged apnea following succinylcholine or mivacurium. Treatment decisions in poisoning must not wait for the laboratory result.

\section*{How This Test Is Performed}
A blood sample is analyzed for enzyme activity. If inherited deficiency is suspected, a dibucaine number or genetic testing may be added to characterize the cause.

\section*{Preparation, Risks, and Limitations}
No fasting is required. Venipuncture is the usual risk. Liver disease, malnutrition, pregnancy, inflammation, and several medicines can alter activity. Results do not reliably quantify the severity of poisoning without exposure history and clinical assessment.

\section*{Understanding Results}
Markedly reduced activity supports significant cholinesterase-inhibitor exposure in the correct setting. Low activity can also be inherited or acquired. A low dibucaine number suggests an atypical enzyme associated with prolonged succinylcholine effect.

\section*{References}
World Health Organization guidance on clinical management of pesticide poisoning.

\clearpage
\section*{Regulatory Status and Authorization Date}
Regulatory status applies to the specific assay, instrument, manufacturer, and intended use rather than to the test name alone. No single approval date can be assigned to this test category. For a commercial assay, verify the current product in the relevant regulator's database: the U.S. Food and Drug Administration (FDA) clearance, approval, or authorization; India's Central Drugs Standard Control Organisation (CDSCO) licence or permission; the European Union In Vitro Diagnostic Regulation (EU IVDR) CE certificate issued through the applicable conformity-assessment route; or the United Kingdom Medicines and Healthcare products Regulatory Agency (MHRA) registration and UKCA/CE status. Laboratory-developed tests may not have an individual FDA, CDSCO, EU IVDR, or MHRA approval date.
\clearpage
